% 题干转写：完成
% 解析转写：完成
% 图形核对：完成
\section{第38届复赛第7题：汽车斜坡加速与车轮附着条件}

\subsection{题干}
如图\ref{fig:cpho38q07}，一简化汽车在倾角为 $\theta$ 的斜坡上运动。前轮（两个前轮合并视为质量 $m$ 的物体）和后轮（两个后轮合并视为质量 $m$ 的物体）均为半径 $r$ 的均质圆柱刚体，前后轮轴相距 $L$；车身质量为 $M$，质心位于两轴正中，距轮轴平面高度为 $h$。轮胎与地面的摩擦系数为 $\mu$，最大静摩擦力按滑动摩擦力处理；忽略空气阻力及车轴摩擦。每个车轮组的转动惯量为
\[
J=\frac12mr^2.
\]

求：（1）无滑动上坡加速时前、后轮正压力 $N_{\text{前}}$、$N_{\text{后}}$ 及前轮不离地的最大安全加速度；（2）发动机向后轮施加顺时针力矩 $\tau_{\text{轮}}$ 时的车身加速度；（3）对给定 A 型电动汽车，分别计算水平路面和 $30^\circ$ 上坡时的最大加速度，并给出不发生滑动所需的最小摩擦系数。给定 $M=1.8\times10^3\,\mathrm{kg}$、$h=0.10\,\mathrm m$、$r=0.40\,\mathrm m$、$L=2.9\,\mathrm m$、$\tau=4.0\times10^2\,\mathrm{N\,m}$、$r_i=10$，车轮质量相对车身质量可忽略，$g=9.8\,\mathrm{m/s^2}$。

\begin{figure}[htbp]
  \centering
  \begin{tikzpicture}[x=.78cm,y=.78cm,>=Stealth,font=\small]
  % Left panel: geometry of the original vehicle-on-slope model.
  \begin{scope}[xshift=-3.9cm,rotate=15]
    \draw[thick] (-1.0,0)--(6.4,0);
    \draw[thick] (-1.0,-.10)--(6.4,-.10);
    \draw[thick,fill=white] (0,.55) circle (.55);
    \draw[thick,fill=white] (5,.55) circle (.55);
    \fill (0,.55) circle (1.4pt) node[below left] {$m$};
    \fill (5,.55) circle (1.4pt) node[below right] {$m$};
    \draw[thick] (-.25,1.15)--(.35,1.65)--(4.65,1.65)--(5.25,1.15)--cycle;
    \fill (2.5,1.38) circle (1.6pt) node[above] {$M$};
    \draw[<->] (0,2.05)--(2.5,2.05) node[midway,above] {$L/2$};
    \draw[<->] (2.5,2.05)--(5,2.05) node[midway,above] {$L/2$};
    \draw[densely dashed] (2.5,1.38)--(2.5,.0);
    \draw[<->] (2.72,.0)--(2.72,1.38) node[midway,right] {$h$};
    \draw[->] (0,.55)--(.48,.55) node[midway,above] {$r$};
    \draw[->,thick] (-.42,.73) arc (145:-170:.45) node[left,xshift=-3pt] {$\tau$};
  \end{scope}
  \draw[thick] (-4.9,-1.42)--(-2.7,-1.42);
  \draw[thick] (-4.9,-1.42)--(-4.9,-.55);
  \draw[->] (-4.9,-1.42) -- (-3.85,-1.14);
  \node[below] at (-3.65,-1.26) {$\theta$};
  \node[align=center] at (-1.9,-2.25) {斜坡汽车几何模型};

  % Right panel: force directions used in the dynamics equations.
  \begin{scope}[xshift=3.3cm,rotate=15]
    \draw[thick] (-1.0,0)--(5.7,0);
    \draw[thick] (-1.0,-.10)--(5.7,-.10);
    \draw[thick,fill=white] (0,.52) circle (.52);
    \draw[thick,fill=white] (4.2,.52) circle (.52);
    \draw[thick] (-.18,1.12)--(.32,1.52)--(3.88,1.52)--(4.38,1.12)--cycle;
    \coordinate (cphoQ07cg) at (2.1,1.3);
    \fill (cphoQ07cg) circle (1.6pt);
    \draw[->,very thick] (2.1,1.3)--(3.25,1.3) node[above] {$a$};
    \draw[->,thick] (0,.03)--(-.62,.03) node[below] {$F_f$};
    \draw[->,thick] (4.2,.03)--(4.98,.03) node[below] {$F_r$};
    \draw[->,thick] (0,.05)--(0,.95) node[left] {$N_f$};
    \draw[->,thick] (4.2,.05)--(4.2,.95) node[right] {$N_r$};
    \draw[->,thick] (-.23,.70) arc (130:-170:.40) node[left] {$\tau$};
  \end{scope}
  \draw[->,thick] (cphoQ07cg)--++(0,-1.08) node[right] {$Mg$};
  \node[align=center] at (5.0,-2.25) {沿坡面动力学示意};
\end{tikzpicture}

  \caption{第38届复赛第7题汽车斜坡模型的TikZ重绘。左图对应原题图7a，右图标示沿坡面动力学中使用的主要作用力和力矩。}
  \label{fig:cpho38q07}
\end{figure}

\subsection{参考解答整理}
\paragraph{（1）正压力与安全加速度。}
取 $x$ 轴沿坡面向上、$y$ 轴垂直坡面向上。前、后轮轴对车身的支撑力满足
\[
F_{x}=M(a+g\sin\theta),\qquad
F_{y,\text{前}}+F_{y,\text{后}}=Mg\cos\theta,
\]
并有相对车身质心的力矩平衡
\[
F_{y,\text{前}}\frac L2+F_xh-F_{y,\text{后}}\frac L2=0.
\]
代入 $F_{y,\text{前}}=N_{\text{前}}-mg\cos\theta$、$F_{y,\text{后}}=N_{\text{后}}-mg\cos\theta$，得
\[
\boxed{N_{\text{前}}=mg\cos\theta+\frac12Mg\cos\theta
-\frac hL M(a+g\sin\theta),}
\]
\[
\boxed{N_{\text{后}}=mg\cos\theta+\frac12Mg\cos\theta
+\frac hL M(a+g\sin\theta).}
\]
令 $N_{\text{前}}=0$，前轮刚好离地，对应安全上限为
\[
\boxed{a_{\text{安全}}=
\frac{(2m+M)L}{2Mh}g\cos\theta-g\sin\theta.}
\]

\paragraph{（2）无滑动驱动时的加速度。}
以逆时针为车轮转动的正方向。无滑动约束给出
\[
\alpha_{\text{前}}=\alpha_{\text{后}}=-\frac ar.
\]
前轮转动方程与 $J=mr^2/2$ 给出
\[
F_{x,\text{前}}r=J\alpha_{\text{前}},
\qquad F_{x,\text{前}}=-\frac m2a.
\]
整车沿斜面方向的运动方程和后轮转动方程为
\[
F_{x,\text{前}}+F_{x,\text{后}}-(M+2m)g\sin\theta=(M+2m)a,
\]
\[
-\tau_{\text{轮}}+F_{x,\text{后}}r=J\alpha_{\text{后}}.
\]
联立可得
\[
\boxed{a=\frac{\tau_{\text{轮}}/r-(M+2m)g\sin\theta}{M+3m}.}
\]

\paragraph{（3）数值结果与附着条件。}
变速器效率为 $100\%$ 时，传递给车轮的力矩为
\[
\tau_{\text{轮}}=\frac{\tau\omega}{\omega_{\text{轮}}}=\tau r_i.
\]
对 A 型车，$m\ll M$，因而
\[
\boxed{a(0^\circ)=5.6\,\mathrm{m/s^2},\qquad
a(30^\circ)=0.66\,\mathrm{m/s^2}.}
\]
前、后轮不打滑应分别满足
\[
|F_{x,\text{前}}|\le\mu N_{\text{前}},\qquad
|F_{x,\text{后}}|\le\mu N_{\text{后}}.
\]
在 $m\ll M$ 的近似下，前轮摩擦力为零。题解代入后轮受力得到
\[
N_{\text{后}}\big|_{\theta=0}=9.2\times10^3\,\mathrm N,
\qquad F_{x,\text{后}}\big|_{\theta=0}=1.0\times10^4\,\mathrm N,
\]
\[
N_{\text{后}}\big|_{\theta=30^\circ}=8.0\times10^3\,\mathrm N,
\qquad F_{x,\text{后}}\big|_{\theta=30^\circ}=1.0\times10^4\,\mathrm N.
\]
因此两种路况分别要求
\[
\boxed{\mu_{\min}(0^\circ)=1.09,\qquad
\mu_{\min}(30^\circ)=1.25.}
\]

\subsection{转写依据}
题干依据 \texttt{CPHO38-Q07\_transcription\_v1.md}；公式依据 \texttt{CPHO38-Q07\_formula\_check\_v1.csv} 与答案第16至18页；题图依据原题图7a重绘。本题第（1）至（3）问评分分别为16分、16分、8分。
